\originalchapter{期末考试}
\psource{8.044，2013 春，Final Exam，五题各 20 分，课程安排考试时长三小时。中文题面；据公开官方答案核对的 AI 独立解答。}
\originalsection{1．二元合金}
\begin{exercise}
$N$ 个 $\alpha$ 原子、$N$ 个 $\beta$ 原子，各有 $N$ 个固定相应位点。零温完美有序，有限温 $n$ 个 $\alpha$ 换至 $\beta$ 位，等量 $\beta$ 占空出的 $\alpha$ 位，失序能 $U=\epsilon n$，$\epsilon>0$ 恒定。
(a) 选择 $n$ 个被 $\beta$ 占据的 $\alpha$ 位有多少方法；(b) 选择 $n$ 个换出的 $\beta$ 位方法数；(c) 求 $S(n)$；(d) 求 $U(T,N)$。
\end{exercise}
\begin{solution}
(a)(b) 各为 $\binom Nn$，两种选择独立且相同种原子不编号，$\Omega=\binom Nn^2$。
(c) $S=2\kb\ln\binom Nn\simeq2\kb[N\ln N-n\ln n-(N-n)\ln(N-n)]$。
(d) 大体系 $1/T=S_n/\epsilon=2\kb\ln[(N-n)/n]/\epsilon$，故 $n=N/[1+e^{\epsilon/(2\kb T)}]$、$U=\epsilon N/[1+e^{\epsilon/(2\kb T)}]$。低温半能隙指数来自两类位点的组态数，不是改写单交换能量 $\epsilon$。
\end{solution}
\originalsection{2．DNA 拉链模型}
\begin{exercise}
$N$ 个碱基对链结闭合能零、打开能 $\epsilon>0$。仅能从一端打开，第 $n$ 结打开前全部 $1,\ldots,n-1$ 必已开，$N$ 很大。
(a) 单链 $Z_1$ 留为和；(b) $\kb T\ll\epsilon$ 时平均开结数；(c) $N\epsilon\gg\kb T\gg\epsilon$ 时用积分近似 $Z_1$；(d) 同温区平均开结数。
\end{exercise}
\begin{solution}
(a) $Z_1=\sum_{n=0}^Ne^{-\beta\epsilon n}$，每个 $n$ 只有一态，不是 $\binom Nn$。
(b) 仅保留 $n=0,1$，$\avg n\simeq e^{-\beta\epsilon}/(1+e^{-\beta\epsilon})\simeq e^{-\beta\epsilon}$。
(c) $Z_1\simeq\int_0^Ne^{-\beta\epsilon n}\dd n=(1-e^{-\beta\epsilon N})/(\beta\epsilon)\simeq\kb T/\epsilon$。
(d) $\avg n=-\epsilon^{-1}\partial_\beta\ln Z_1\simeq\kb T/\epsilon$。这是远未全开的中间“高温”区；真正 $\kb T\gg N\epsilon$ 时全部 $N+1$ 态等概率，平均 $N/2$。
\end{solution}
\originalsection{3．自旋波}
\begin{exercise}
低温铁磁体的相关自旋偏转为正常模，每波矢一个谐振子。三维周期边界为\[\vct k=(2\pi m/L_x,2\pi n/L_y,2\pi l/L_z),\]整数遍取，低温立方晶格近似 $\omega=ak^2$。
(a) 波矢密度；(b) 求并画 $D(\omega)$；(c) 热容留无量纲积分，可从开始略去不随温度变的零点能；(d) 是否有能隙行为并解释。
\end{exercise}
\begin{solution}
(a) $D_k=V/(2\pi)^3$。(b) 球壳换元，$D(\omega)=V\sqrt\omega/(4\pi^2a^{3/2})$，平方根从零上升。
(c) Bose 模式热能求导得
\[C_V=\frac{V\kb}{4\pi^2}\left(\frac{\kb T}{\hbar a}\right)^{3/2}\int_0^\infty\frac{x^{5/2}e^x}{(e^x-1)^2}\dd x\propto T^{3/2}.\]
(d) 连续体相中 $k\to0$ 无间隙，故幂律而非激活指数；有限尺寸及磁各向异性可产生额外截止或间隙，题设忽略它们。
\end{solution}
\begin{center}\begin{tikzpicture}\begin{axis}[width=7cm,height=3.4cm,xmin=0,xmax=3,ymin=0,ymax=1.8,xlabel={归一化 $\omega$},ylabel={归一化 $D$}]\addplot[domain=0:3,thick]{sqrt(x)};\end{axis}\end{tikzpicture}\end{center}
\originalsection{4．石墨烯的线性色散}
\begin{exercise}
单层碳片导带电子二维色散 $\epsilon=\hbar v|\vct k|$，样品 $L_x,L_y$ 周期边界，$\vct k=(2\pi m/L_x,2\pi n/L_y)$。
(a) 波矢密度；(b) 每电子自旋简并 2，求导带 $D_c(\epsilon)$；(c) 价带负能 DOS 为镜像，$D_v(\epsilon)=D_c(-\epsilon)$，两带无隙相触，画总 DOS；(d) 零温价带全满导带全空，求化学势与升温时移动方向；(e) 以 $U(0)=0$ 为能量基准，求激发内能，无量纲积分不求值，注意从负能态至对应正能态耗双倍能；(f) 求低温 $C_A\propto T^n$ 的指数。
\end{exercise}
\begin{solution}
(a) $D_k=A/(2\pi)^2$。(b) 包含题设自旋二重，$D_c=A\epsilon/(\pi\hbar^2v^2)$（$\epsilon>0$）。真实石墨烯还有谷简并，本题未给，不能在作答中暗加因子。
(c) $D(\epsilon)=A|\epsilon|/(\pi\hbar^2v^2)$，V 字形。(d) 粒子--空穴完全对称且总数中性，$\mu=0$ 并随温度保持零。
(e) 导带电子与价带空穴各贡献一次，
\[U=\frac{2A(\kb T)^3}{\pi\hbar^2v^2}\int_0^\infty\frac{x^2}{e^x+1}\dd x.\]
(f) $C_A\propto T^2$，$n=2$。带宽以外此无限线性色散模型失效。
\end{solution}
\begin{center}\begin{tikzpicture}\begin{axis}[width=7cm,height=3.4cm,xmin=-2,xmax=2,ymin=0,ymax=2.1,xlabel={归一化 $\epsilon$},ylabel={归一化 $D$}]\addplot[domain=-2:2,thick]{abs(x)};\end{axis}\end{tikzpicture}\end{center}
\originalsection{5．凝聚与加粒子的内能导数}
\begin{exercise}
三维无自旋理想玻色原子 $D(\epsilon)=V(2m/\hbar^2)^{3/2}\sqrt\epsilon/(4\pi^2)$。
(a) 推导 $T,n=N/V,\mu$ 一般关系，以 $y=\beta\mu$ 引入无量纲积分 $I(y)$，不求值；(b) 求 BEC 起始密度温度关系，不求积分；(c) 用 Maxwell 关系证明 $(\partial U/\partial N)_{T,V}=\mu-T(\partial\mu/\partial T)_{N,V}=(\partial y/\partial\beta)_{N,V}$，并说明凝聚相此导数为零；(d) 无凝聚时用 (a) 证明该导数不会为负。
\end{exercise}
\begin{solution}
令 $C=(2m/\hbar^2)^{3/2}/(4\pi^2)$。(a) 非凝聚连续谱满足 $n=C(\kb T)^{3/2}I(y)$，
\[I(y)=\int_0^\infty\frac{\sqrt x}{e^{x-y}-1}\dd x,\qquad y<0.\]
一般相还应加基态 $n_0$；不能让无基态的积分超过它的最大容量。
(b) 起始 $n_0=0,y\to0$，$n=C(\kb T_c)^{3/2}I(0)$，即 $T_c\propto n^{2/3}$。
(c) $U=F+TS$，$F_N=\mu$，$S_N=-\mu_T$，故等温加粒子导数为所求。$y=\beta\mu$ 的链式法则给第二等号。凝聚热力学极限 $\mu=0$、激发云只依赖 $T,V$，加粒子进入零能基态，不增内能；有限体系存在小修正。
(d) 固定 $n$ 对 $n=C\beta^{-3/2}I(y)$ 求 $\beta$ 导数，$y_\beta=3I/(2\beta I')>0$，其中 $I'>0$ 来自逐项正导数，且 $y<0$ 可控积分收敛。接近临界 $I'\to\infty$ 时导数趋零，连续衔接凝聚侧。
\end{solution}
